For every $m\in M$, \exref{4.21} gives
\[r(N:m)=\bigcap_{i=1}^nr(Q_i:m)=\bigcap_{m\notin Q_i}\mathfrak{p}_i.\]
If this ideal is prime, it equals one of the $\mathfrak{p}_i$, since a prime containing a finite intersection of prime ideals contains one of them. Conversely, irredundancy gives $m\in\bigcap_{j\ne i}Q_j\setminus Q_i$, for which $r(N:m)=\mathfrak{p}_i$. Thus the $\mathfrak{p}_i$ are precisely the prime ideals of the form $r(N:m)$ and depend only on $N\subseteq M$.

Since $(N:m)=\operatorname{Ann}(m+N)$ and the $Q_i/N$ form a minimal primary decomposition of zero in $M/N$, the same primes belong to that zero submodule.
